% Abhakseite „Das kann ich" – Zone nur im Gesamt (\mitzone)
\newcommand{\abhakverf}[1]{\abhakgruppe{\normalfont\itshape #1}}
\begin{abhakseite}
\ifdefined\mitzone
\abhakgruppe{Kennst du schon}
\abhak{1}{Ich kann ablesen, wie viel Prozent eines Streifens gefärbt sind.}
\abhak{2}{Ich kann einen Anteil in Prozent in einen Streifen einzeichnen.}
\abhak{3}{Ich kann von 10 \% auf mehr Prozent und auf das Ganze schließen.}
\abhak{4}{Ich kann einen Anteil als Bruch schreiben und auf den Nenner 100 erweitern.}
\abhak{5}{Ich kann einen Bruch als Dezimalzahl schreiben und umgekehrt.}
\abhak{6}{Ich kann auf eine Stelle nach dem Komma runden.}
\abhak{7}{Ich kann durch 100 teilen und mit einer Dezimalzahl malnehmen.}
\abhak{8}{Ich kann mit einer Tabelle über 1 Stück rechnen.}
\abhak{9}{Ich kann einen Bruchteil einer Größe ausrechnen.}
\abhak{10}{Ich finde den Fehler beim Umwandeln in eine Dezimalzahl.}
\abhak{11}{Ich kann Hundertstel als Dezimalzahl schreiben und umgekehrt.}
\fi
\abhakgruppe{Einheit 1 · Prozentsatz berechnen}
\abhak{12}{Ich erkenne, was das Ganze ist.}
\abhak{13}{Ich kann einen Streifen in gleich große Teile einteilen.}
\abhakverf{Anteil in Prozent am Streifen}
\abhak{14}{Ich kann am Streifen ablesen, wie viel Prozent ein Teil vom Ganzen ist.}
\abhakverf{Anteil in Prozent durch Erweitern}
\abhak{15}{Ich kann durch Erweitern ausrechnen, wie viel Prozent ein Teil vom Ganzen ist.}
\abhak{16}{Ich kann durch Erweitern ausrechnen, wie viel Prozent ein Teil vom Ganzen ist – weiter.}
\abhakverf{Anteil in Prozent mit dem Taschenrechner}
\abhak{17}{Ich kann mit dem Taschenrechner ausrechnen, wie viel Prozent ein Teil vom Ganzen ist.}
\abhakverf{Anteil in Prozent aus einem Sachtext}
\abhak{18}{Ich kann in einem Sachtext ausrechnen, wie viel Prozent ein Teil vom Ganzen ist.}
\abhak{19}{Ich kann in einem Sachtext ausrechnen, wie viel Prozent ein Teil vom Ganzen ist – weiter.}
\abhak{20}{Ich kann zu einem Prozentsatz passende Zahlen finden.}
\abhak{21}{Ich finde den Fehler, wenn jemand ausrechnet, wie viel Prozent ein Teil ist.}
\abhak{22}{Ich kann begründen, warum man durch das Ganze teilt.}
\abhak{23}{Ich kann mit Prozent vergleichen, wer besser trifft.}
\abhakgruppe{Einheit 2 · Prozentwert berechnen}
\abhakverf{Teil am Streifen}
\abhak{24}{Ich kann am Streifen ausrechnen, wie viel ein einfacher Anteil in Prozent von einer Größe ist.}
\abhak{25}{Ich kann am Streifen ausrechnen, wie viel ein einfacher Anteil in Prozent von einer Größe ist – weiter.}
\abhakverf{Teil über 1 \% und über 10 \%}
\abhak{26}{Ich kann über 1 \% ausrechnen, wie viel ein Anteil in Prozent von einer Größe ist.}
\abhak{27}{Ich kann über 10 \% ausrechnen, wie viel ein Anteil in Prozent von einer Größe ist.}
\abhakverf{Teil mit der Dezimalzahl}
\abhak{28}{Ich kann einen Anteil in Prozent als Dezimalzahl mit einer Größe malnehmen.}
\abhakverf{Teil im Sachtext: sparen und bezahlen}
\abhak{29}{Ich kann ausrechnen, wie viel man spart und wie viel man noch bezahlt.}
\abhak{30}{Ich kann ausrechnen, wie viel mehr als 100 \% von einer Größe sind (je nach Buch erst in Klasse 8).}
\abhak{31}{Ich finde den Fehler, wenn jemand einen Anteil in Prozent von einer Größe ausrechnet.}
\abhak{32}{Ich kann begründen, warum man für 1 \% durch 100 teilt.}
\abhak{33}{Ich kann mit der Mehrwertsteuer Preise vergleichen.}
\abhakgruppe{Einheit 3 · Grundwert berechnen}
\abhak{34}{Ich erkenne, was gegeben und was gesucht ist.}
\abhakverf{Das Ganze am Streifen}
\abhak{35}{Ich kann am Streifen ausrechnen, wie viel das Ganze ist.}
\abhakverf{Das Ganze über 1 \%}
\abhak{36}{Ich kann über 1 \% ausrechnen, wie viel das Ganze ist.}
\abhakverf{Das Ganze im Sachtext}
\abhak{37}{Ich kann in einem Sachtext ausrechnen, wie viel das Ganze ist.}
\abhak{38}{Ich kann unterscheiden, ob der Teil, der Anteil in Prozent oder das Ganze gesucht ist.}
\abhak{39}{Ich finde den Fehler, wenn jemand das Ganze ausrechnet.}
\abhak{40}{Ich kann begründen, warum man für das Ganze mal 100 rechnet.}
\abhak{41}{Ich kann mit dem Ganzen planen, wie viele Plätze gebraucht werden.}
\abhakgruppe{Einheit 4 · Prozentuale Veränderung}
\abhak{42}{Ich erkenne, ob etwas um einen Prozentsatz oder auf einen Prozentsatz geändert wird.}
\abhak{43}{Ich erkenne, welcher Wert der alte ist.}
\abhakverf{Neuer Wert: Teil dazu oder weg}
\abhak{44}{Ich kann ausrechnen, wie groß ein Wert nach einer Erhöhung oder Senkung ist.}
\abhak{45}{Ich kann ausrechnen, wie groß ein Wert nach einer Erhöhung oder Senkung ist – weiter.}
\abhakverf{Neuer Wert mit dem Faktor}
\abhak{46}{Ich kann den neuen Wert in einem Schritt ausrechnen (kommt in Klasse 9; am Gymnasium schon in Klasse 7).}
\abhakverf{Veränderung in Prozent}
\abhak{47}{Ich kann ausrechnen, um wie viel Prozent sich ein Wert verändert hat.}
\abhakverf{Alter Wert aus neuem Wert}
\abhak{48}{Ich kann aus dem neuen Wert den alten ausrechnen.}
\abhakverf{Brutto und Netto}
\abhak{49}{Ich kann Preise mit und ohne Mehrwertsteuer ausrechnen (neu in diesem Jahr).}
\abhakverf{Prozentpunkte}
\abhak{50}{Ich kann Prozentpunkte und Prozent unterscheiden.}
\abhakverf{Steigung in Prozent}
\abhak{51}{Ich kann eine Steigung in Prozent ausrechnen (kommt in Klasse 10).}
\abhak{52}{Ich finde den Fehler, wenn jemand eine Veränderung in Prozent ausrechnet.}
\abhak{53}{Ich kann begründen, welcher Wert 100 \% ist.}
\abhak{54}{Ich kann zwei Rabattangebote vergleichen.}
\end{abhakseite}
